\section{Introduction}
\label{sec:introduction}

Let $N=2^L$ with $L\ge2$ and associate a real vector $x\in\R^N$ with the residue class
\begin{equation}
  x(z)=\sum_{n=0}^{N-1}x_nz^n
  \quad\text{in}\quad
  R_N=\R[z]/(z^N-1).
  \label{eq:intro-polynomial}
\end{equation}
The discrete Fourier transform is the simultaneous evaluation of this class at
the $N$ roots of unity.  Complex linear algebra is therefore a coordinate
choice, not a requirement of the polynomial problem: conjugate roots have the
real minimal polynomial $z^2-2\cos\theta\,z+1$.

Bruun organized these real factors into a sparse recursive reduction
\cite{bruun1978}.  The construction was remarkable for exposing a real
coefficient route to the spectrum, but three questions remain mathematically
prior to any implementation question.  What is the local inverse of a
trinomial split?  Which coordinates make the factorization metrically regular?
How does the native leaf order identify the ordinary Fourier frequencies?
Answering only the first question gives an exact polynomial reconstruction;
answering all three identifies a stable Fourier operator.

The present investigation began from those questions.  It then asked whether
alternative factor bases changed the local multiplication geometry; whether
blocks of equal-angle nodes possessed a common real representation; which
spectral orders preserve subtree locality; and whether the same factor tree can
be traversed upward rather than downward.  Those questions lead, respectively,
to a normalized quadratic chart, an isoclinic block form, a classification of
admissible leaf orders, and a decimation-in-time dual.  Two further questions
extend the algebra.  Replacing roots of $z^N-1$ by roots of $z^N+1$ produces an
odd-frequency transform, while shearing the intermediate coordinates produces
a diagonal phase-packet transform.  A final question asks whether the signed
linear cells admit a conservative nonnegative or exact equilibrium
realization.

The principal results are as follows.
\begin{enumerate}
  \item An explicit sibling merge gives a two-sided Chinese-remainder inverse
        for every binomial and trinomial node.
  \item The normalized terminal coordinate $e_\theta$ squares to $-1$ and
        converts evaluation into the standard real and imaginary Fourier
        coordinates.
  \item Every normalized node is $\sqrt{2}$ times an orthogonal map, giving
        $B_N^{\trans}B_N=NI$ and condition number one.
  \item A commuting reduction-evaluation diagram proves
        $P_NB_N=\DFT_N^{\R}$ exactly; the factorization is not an approximation
        to the DFT.
  \item DIF, DIT, and diagonal phase-packet schedules are boundary-permuted
        factorizations of that same operator.
  \item The roots of $z^N+1$ give an exact real two-plane odd-frequency
        transform with the same orthogonality law.
  \item A positive-cone intertwining identity gives exact signed projection
        through arbitrary compositions, and diagonal gauges convert the cells
        into conservative transports or convex equilibrium maps.
\end{enumerate}

The mathematical development is intentionally ordered by dependency.  Section
\ref{sec:bruun} reconstructs Bruun's factor tree and the questions that resume
its thread.  Sections \ref{sec:crt}--\ref{sec:fourier-identity} give the inverse,
normalization, and exact Fourier identity.  Sections
\ref{sec:isoclinic}--\ref{sec:dip} develop the block, ordering, DIT, ODFT, and
DIP consequences.  Section \ref{sec:physical} proves physical realizability in
an explicit ideal linear category.  Exact examples and finite-precision audits
follow only after the algebraic results.
